% s-077-p077-m002; original PDF 77
\Aphorism{s-077-p077-m002}{67}
Wie sehr in der Wissenschaft ausserwissenschaftliche Neigungen und Abneigungen mitsprechen, sieht man daraus, dass in der typischen wissenschaftlichen Fehde stets einem „Das ist“ ein „Das kann nicht sein“ gegenübergestellt wird; man streitet nicht um wirklich oder unwirklich, möglich oder unmöglich, sondern sofort um die Extreme wirklich oder unmöglich. Diese Verschärfung der Alternative, wo Jeder durchaus mehr widerlegen will als der Andere behauptet, deutet auf ein Für und Wider ausserhalb der reinen Erkenntniss; solchen Staub wirbelt kein Kampf unschuldiger Meinungen auf, sondern höchstens ein Kampf der Interessen, Personen, „Ueberzeugungen“, Glaubensrichtungen. Wir kennen den Ton, der jede Skepsis, jedes vorsichtige Offenhalten und Vertagen der Entscheidung mit Heftigkeit zurückweist; aber es ist nicht der Ton der Wissenschaft.
\par

